\section{Operational View Specification}

\subsection{System Performance}

\subsubsection{Response Time}
The system operates in near real-time. Motion events trigger lighting and ventilation in approximately 1--2 seconds. In the wudu area, abnormal continuous flow triggers valve closure in approximately 2--3 seconds to reduce waste while limiting false positives.

\begin{figure}[htbp]
  \centering
  \includegraphics[width=0.7\textwidth]{8.1.1.jpg}
  \caption{Operational response-time trend}
\end{figure}

Air quality is monitored periodically and ventilation is adjusted when thresholds are exceeded. Temperature and humidity are also tracked to maintain comfort. The system handles simultaneous sensor inputs during peak times without noticeable degradation.

The graph proves that, considering the mosque's system operates over a local network, the typical response time for light activation or ventilation devices will not exceed 2 seconds.

The system continuously monitors CO$_2$ and other pollutants and responds when levels exceed safe ventilation thresholds. Sensors provide regular readings, enabling quick response to environmental changes while maintaining indoor comfort.

\subsubsection{Data Update Frequency}
\begin{itemize}
  \item Occupancy: real-time update per entry/exit event.
  \item Air quality: every 1 minute (balanced responsiveness and stability).
  \item Water usage: continuously monitored in real time.
\end{itemize}

The system updates occupancy immediately after every entry or exit event so the dashboard always reflects the current number of worshippers inside the prayer hall. This immediate response is important for crowd management during high-attendance periods.

When the mosque becomes crowded, ventilation can also be triggered proactively based on occupancy levels before the next air-quality update cycle to maintain airflow and comfort.

For water usage, the system continuously tracks total consumption and detects abnormal flow behavior such as leaks or taps left open. If abnormal flow is detected, the system can trigger alerts and automatically close the smart valve to reduce waste.

\subsubsection{Throughput}
The system processes motion, water-flow, and air-quality streams in parallel, enabling simultaneous monitoring and decision-making during high-demand periods such as Friday prayer and Ramadan.

The system is designed to handle multiple sensors simultaneously by collecting occupancy, water, and air-quality data in parallel and forwarding readings to the controller without blocking other operations. This ensures the dashboard remains responsive even during peak activity periods.

\subsubsection{Efficiency}
Lights and fans operate only when needed based on occupancy and schedule, and water control logic closes supply during leak-like behavior to reduce waste.

The system avoids unnecessary device activation by ensuring that lights and ventilation operate only when occupancy is detected and prayer-time conditions require activation. Water efficiency is also improved through continuous flow monitoring and automatic shutoff during abnormal usage.

\subsection{System Reliability}

\subsubsection{Sensor Failure Handling}
Fallback mechanisms keep operation safe during sensor faults.

\begin{itemize}
  \item Motion sensor failure: admin manual override controls lighting and fans.
  \item Flow sensor failure: manual valve override + admin alert.
  \item Air-quality sensor failure: ventilation follows safe default schedule.
\end{itemize}

The system remains operational even if one sensor becomes unavailable, allowing maintenance to occur without interrupting mosque services.

\subsubsection{Network Failure}
Core control continues locally if network/server connectivity is lost. Dashboard synchronization resumes automatically after restoration.

Essential functions such as lights, fans, and water shutoff continue operating locally through the controller even during temporary internet or server outages.

\subsubsection{Power Failure}
On recovery, the system returns to predefined safe startup states. Critical monitoring may be maintained through optional backup power.

If backup power is available, critical monitoring functions can continue during outages. After power restoration, devices return to predefined safe states instead of restarting unpredictably.

\subsection{System Scalability}

\subsubsection{More Sensors}
Additional occupancy, flow, and IAQ sensors can be integrated with the same controller logic.

New sensors can be added gradually as mosque requirements expand. Additional occupancy sensors, flow sensors, and air-quality sensors can extend monitoring coverage without redesigning the entire system.

\subsubsection{Spreading to Other Mosques}
The architecture can be replicated across multiple mosques, each with local operation and data separation.

The same framework can be reused in multiple mosque deployments while maintaining independent local operation and monitoring for each location.

\subsubsection{Modular Architecture}
Occupancy, water, air-quality, and dashboard modules can evolve independently, reducing maintenance risk and easing upgrades.

The window-ventilation and humidity-comfort modules operate independently and were added without modifying the core occupancy or air-quality logic, demonstrating the modular design principle in practice.

This modular approach allows updates, repairs, and future feature additions to occur with minimal disruption to other system components.

\subsection{System Security}

\subsubsection{Network Protection}
\begin{itemize}
  \item Segmented network for IoT devices
  \item Restricted dashboard/control access
  \item Firewall/gateway policy enforcement
  \item Monitoring for suspicious behavior and repeated login attempts
\end{itemize}

All IoT devices communicate through a secure local network rather than being openly exposed. This reduces external interference risks and improves protection for sensors and actuators.

\begin{itemize}
  \item Sensors, controllers, and actuators operate on a separate secured network segment.
  \item Only approved users and services can access dashboard and device controls.
  \item Firewall and gateway policies reduce unauthorized access risks.
  \item Network traffic is monitored for suspicious activity and repeated login attempts.
\end{itemize}

\subsubsection{Access Control}
Role-based permissions, authenticated overrides, and action logging are applied for accountability.

Only authorized users can access the dashboard or override lighting, ventilation, and water-control systems. The administrator receives full control privileges while other users may receive limited or read-only access.

Manual overrides require authentication credentials, and all override actions are logged for accountability and auditing purposes.

Authentication flow:
\begin{center}
User $\rightarrow$ Login Page $\rightarrow$ Authentication Check $\rightarrow$ Access Granted $\rightarrow$ Dashboard
\end{center}

\subsubsection{Data Encryption}
WPA2-PSK secures Wi-Fi communication between sensors, controllers, and dashboard. AES encryption protects traffic confidentiality and integrity.

Data transmitted across the wireless network is protected using WPA2-PSK encryption. AES encryption secures data confidentiality and helps prevent unauthorized interception or modification of occupancy, water-usage, and device-control data.

\begin{itemize}
  \item Secures communication between sensors and controllers.
  \item Protects local dashboard traffic.
  \item Restricts network access to authorized devices only.
  \item Provides strong security while remaining suitable for embedded IoT devices.
\end{itemize}

\subsubsection{Device Filtering}
MAC allow-list filtering blocks unknown devices from joining. It is used together with WPA2-PSK and access control (not standalone security).

\begin{itemize}
  \item Only trusted devices are allowed to join the network.
  \item Unknown devices are denied access.
  \item Suitable for small controlled IoT deployments.
  \item Used together with WPA2-PSK and role-based access control.
  \item Simplifies secure enrollment of new sensors and controllers.
\end{itemize}

\subsubsection{Privacy Protection}
The system stores aggregate occupancy counts only and avoids personal identification. PIR-based sensing does not capture images or personally identifiable information.

\subsection{Energy Efficiency}

\subsubsection{Prayer Time Control}
Prayer-time-aware control avoids continuous operation outside worship periods.

Prayer-time automation ensures that lighting and ventilation activate only during scheduled worship windows, reducing unnecessary energy consumption during low-use periods.

\subsubsection{Occupancy-Based Activation}
Devices activate only when people are present, reducing unnecessary consumption.

Occupancy detection prevents unnecessary activation when the prayer hall is empty, even if the prayer-time window is active.

\subsubsection{Automatic Shutdown}
Lights and fans turn off after inactivity timeout, minimizing waste.

Automatic shutdown after inactivity reduces unnecessary energy usage when worshippers leave or when only part of the mosque is occupied.

\subsection{Safety Considerations}

\subsubsection{Device Overheating}
Thermal-safe operation requires ventilation, safe power design, and alerts/shutdown under abnormal heating.

The system includes safeguards to prevent overheating of controllers, relays, and power modules during long operating periods.

\subsubsection{Water Valve Reliability}
Valve closure is validated over a short window to avoid false shutdown during normal use interruptions.

The water valve verifies abnormal flow conditions for a short duration before activation, reducing false shutdowns during normal ablution usage.